\section{SF-E849: fixed-ideal deformations exclude non-CI point modifications}\label{SF-E849:section}

Author: /root/sixfolds, GPT-6 Astra Pro; inherited effort not exposed.
Date: 2026-11-22. Author, root and independent whole-proof readings passed without mathematical repair. The frozen mathematical original is retained.
This canonical conversion has a separate full-text check.
This extends the jointly developed E848 method by keeping the arbitrary
quotient ideal fixed. No lifting of an extension class along a deformation
of that ideal is assumed.

All geometric conclusions below concern the auxiliary abelian-fourfold
object problem. They do not assert a transfer to an imaginary-quadratic
Weil eightfold or solve the original target-eightfold seed problem.

\subsection{1. Statements}\label{SF-E849:s1}

Let \(R=\mathbf C[[x_1,x_2,x_3,x_4]]\), and let \(I\) be any proper finite-colength ideal.
Set \(A=R/I\), and define

\[
 \ell=\dim_{\mathbf C}A,\qquad
 a=\dim_{\mathbf C}\operatorname{Hom}_R(I,A),\qquad
 \delta=a-3\ell.
 \tag{NI.1}
\]

There is no complete-intersection, Gorenstein, reducedness, or grading
hypothesis on \(A\). For \(q\ge 1\), consider any actual extension

\[
 0\longrightarrow A^q\longrightarrow M_b\longrightarrow I^q
       \longrightarrow0,\qquad
 b\in\operatorname{Mat}_q(V),\quad
 V=\operatorname{Ext}^1_R(I,A).
 \tag{NI.2}
\]

Here \(\operatorname{Mat}_q(V)\) denotes \(q\)-by-\(q\) matrices with entries in the finite
A-module \(V\). No choice of a free A-basis of \(V\) is made or required.
The extension class \(b\) may be zero; no simultaneous scalar-block
form is assumed.

\textbf{Local theorem SF-E849a.} For every (NI.2),

\[
 \dim_{\mathbf C}\operatorname{Ext}^2_R(M_b,M_b)
                  \ \ge\ 2q^2\delta.
 \tag{NI.3}
\]

The inequality is useful when \(\delta\) is positive; no positivity is
asserted for arbitrary \(I\).

For the global statement, let X be a complex abelian fourfold and \(E\)
a coherent sheaf whose canonical map to its double dual has locally
free hull and finite kernel and cokernel:

\[
 0\longrightarrow T\longrightarrow E\longrightarrow
 H=E^{\vee\vee}\longrightarrow Q\longrightarrow0.
 \tag{NI.4}
\]

Suppose that at a point \(p\) its actual completed module is

\[
 \widehat E_p\simeq M_b\oplus R^s,\qquad
 b\in\operatorname{Mat}_m(V),\quad m>0,\quad s\ge0,
 \qquad \delta\ge4.
 \tag{NI.5}
\]

Thus its finite torsion and hull quotient at \(p\) are both free modules
of rank m over one common algebra \(A\), which need not be a complete
intersection.

\textbf{Global conditional theorem SF-E849b.} Fix any specified effective
coherent finite translation/tensor action \(G\) in the supplied S21
framework, together with an actual coherent equivariant structure on \(E\).
If a class \(\xi\) in \(HH^2(X)\) has nonzero holomorphic bivector
component and its actual ordinary evaluation \(\rho_E(\xi)\) is zero,
then the full semiregularity map on invariant Ext has kernel of
dimension at least \(2\delta-6\), hence at least two.

\textbf{Balanced-RM consequence SF-E849c.} Under the genuine RM tangent
splitting and compatible polarization \(\theta=\theta_++\theta_-\), a sheaf
satisfying (NI.4)--(NI.5) and

\[
 e^B\operatorname{ch}(E)
   =r(1-\nu\theta_+^2/2)(1-\nu\theta_-^2/2),
 \qquad r>0,\quad\nu>0,
 \tag{NI.6}
\]

for an algebraic rational divisor B cannot be weakly \(G\)-semiregular
for any such specified action. For the identity group alone, the
local theorem already proves failure of semiregularity without the
character (NI.6) or the evaluation hypothesis.

The strict non-CI example is \(A=R/\mathfrak m^2\). We will prove

\[
       \ell=5,\qquad a=40,\qquad \delta=25,
       \qquad \operatorname{ext}^2_R(M_b,M_b)\ge50q^2.
 \tag{NI.7}
\]

Its socle has dimension four, so it is not a complete intersection.
This genuinely adds a local algebra not covered by E848.

\subsection{2. Elementary duals and the finite algebra's self-Ext dimensions}\label{SF-E849:s2}

Every finite-length \(R\)-module has \(\operatorname{Ext}^j(-,R)=0\) for \(j<4\). This follows
from a composition series and the coordinate Koszul resolution of
the residue field. Applying these vanishings to \(0\to I\to R\to A\to 0\) gives

\[
 \operatorname{Hom}_R(I,R)=R,\qquad
 \operatorname{Hom}_R(I,I)=R,\qquad
 \operatorname{Hom}_R(M_b,R)=R^q,\qquad
 M_b^{**}=R^q.
 \tag{NI.8}
\]

For the second equality, compose an endomorphism of \(I\) with \(I\to R\).
The first equality makes it multiplication by an element of \(R\),
and every such multiplication preserves the ideal \(I\). The other
equalities are induced by the actual quotient and coordinate hull
maps, not just abstract isomorphisms. The hull image is \(I^q\), with
kernel and cokernel both \(A^q\).

The same exact sequence, now with target \(A\), identifies
\(\operatorname{Ext}^1_R(A,A)=\operatorname{Hom}_R(I,A)\). Indeed \(\operatorname{Hom}_R(A,A)\to \operatorname{Hom}_R(R,A)\) is an
isomorphism, both identified with \(A\), and restriction from \(R\) to \(I\)
is zero. It also identifies \(V=\operatorname{Ext}^2_R(A,A)\).

For finite-length perfect complexes on the regular local ring \(R\),
local duality with normalized dualizing complex \(R[4]\) and tensor-Hom
adjunction give perfect pairings

\[
 \operatorname{Ext}^i_R(A,A)^\vee
       \simeq\operatorname{Ext}^{4-i}_R(A,A).
 \tag{NI.9}
\]

The dual here is the complex-vector-space dual after a coordinate
volume fixes the residue identification. The supported complex
\(\operatorname{RHom}_R(A,A)\) has Euler characteristic zero: apply \(\operatorname{Hom}_R(-,A)\) to
a finite free resolution of \(A\); its alternating ranks sum to
\(\operatorname{rank}_R(A)=0\). Since \(\operatorname{End}_R(A)=A\), (NI.9) and this Euler identity give

\[
 \bigl(\operatorname{ext}^i_R(A,A)\bigr)_{i=0}^4
        =(\ell,\ a,\ 2a-2\ell,\ a,\ \ell),
 \qquad \dim_{\mathbf C}V=2a-2\ell.
 \tag{NI.10}
\]

In particular this calculation is not an assumption that \(I\) has
four generators or that \(V\) is a free A-module. It only uses finite
support, regularity of \(R\), and the displayed elementary Hom
identifications.

The precise published duality inputs are those already checked in
E847--E848: Stacks Section 47.18, tag 0A81, Theorem 47.18.3 and
Lemma 47.18.4,\cite{SFStacksE847LocalDuality} and the residue identification of Lipman Lemma 7.3/
Theorem 7.4, pp.~64--69.\cite{SFLipman1984} No duality is applied to an endomorphism
complex of positive rank without its finite-support cone.

\subsection{3. The endomorphism image and the exact fixed-ideal tangent kernel}\label{SF-E849:s3}

For \(M=M_b\) define

\[
 \begin{split}
 K&=\operatorname{Hom}_R(M,A^q),&
 k&=\dim_{\mathbf C}K,\\
 a_L&=\dim_{\mathbf C}\{t\in\operatorname{Mat}_q(A):tb=0\},\\
 \mathcal B&=\operatorname{im}\bigl(
       \operatorname{End}_R(M)\longrightarrow\operatorname{Mat}_q(R)
                                    \bigr),&
 c&=\operatorname{length}(\operatorname{Mat}_q(R)/\mathcal B).
 \end{split}
 \tag{NI.11}
\]

The map to the matrix ring is the map on the canonical hull in
(NI.8). Its kernel is \(K\). Applying \(\operatorname{Hom}_R(-,A^q)\) to (NI.2) gives
\(k=q^2a+a_L\). Every homomorphism of \(I^q\) is the restriction of a
unique matrix over \(R\), by \(\operatorname{Hom}_R(I,I)=R\).

The Yoneda criterion for an endomorphism of the extension is the
equality of its pushforward and pullback classes: an endomorphism
\(t\) of \(A^q\) and a hull matrix \(r\) lift precisely when \(tb=b\bar r\).
Scalar multiplication on the two arguments of Ext agrees, so
this class depends only on \(r\) modulo \(I\). Equivalently \(V\) is
annihilated by \(I\), since its target \(A\) is. Consequently
\(\operatorname{Mat}_q(I)\) is contained in \(\mathcal B\), and

\[
 \begin{split}
 \overline{\mathcal B}
  &=\{r\in\operatorname{Mat}_q(A):
                   \text{there exists }t,\ tb=br\},\\
 \dim\overline{\mathcal B}&=q^2\ell-c,\\
 \psi:\operatorname{Mat}_q(A)^2&\longrightarrow
                 \operatorname{Mat}_q(V),\qquad
                    \psi(t,r)=tb-br,\\
 \dim\operatorname{im}\psi&=q^2\ell+c-a_L.
 \end{split}
 \tag{NI.12}
\]

The last line follows by projecting \(\ker\psi\) to its second factor:
the image is \(\overline{\mathcal B}\) and the kernel has dimension \(a_L\). These are
ordinary finite-dimensional complex-vector-space dimensions.

Now keep \(I\) and \(A\) fixed. Let \(D=\mathbf C[\varepsilon]/(\varepsilon^2)\), \(S=R\otimes_{\mathbf C}D\),
\(I_D=I\otimes_{\mathbf C}D\), and \(A_D=A\otimes_{\mathbf C}D\). A finite free \(R\)-resolution
tensored with \(D\) proves

\[
 \operatorname{Ext}^1_S(I_D^q,A_D^q)
       =\operatorname{Mat}_q(V)\otimes_{\mathbf C}D.
 \tag{NI.13}
\]

Thus each \(b+\varepsilon\beta\) defines an actual extension \(M_{\varepsilon}\)
of the fixed \(D\)-flat modules \(I_D^q\) and \(A_D^q\), and hence a \(D\)-flat
module deformation of \(M\). No deformation of the ideal has been
introduced, and there is no obstruction to lifting \(b\) along a
varying quotient to be discarded.

Tensoring the finite free resolution also gives
\(\operatorname{Ext}^j_S(A_D,S)=0\) for \(j<4\). As in (NI.8), this proves the actual
relative hull identification \(M_\varepsilon^{**}=S^q\), with image \(I_D^q\)
and kernel \(A_D^q\), compatibly with reduction. This proof works
over the nonreduced \(S\) and uses no general double-dual base-change
assertion.

The family defines a linear tangent map
\(L:\operatorname{Mat}_q(V)\to \operatorname{Ext}^1_R(M,M)\). A zero tangent class is exactly a
trivial \(S\)-module deformation with identity special fiber: a
splitting of \(0\to M\to M_{\varepsilon}\to M\to 0\) yields such a trivialization,
and conversely. It induces automorphisms \(1+\varepsilon t\) on the
canonical kernel \(A_D^q\) and \(1+\varepsilon r\) on the canonical hull.
Every hull matrix preserves \(I_D^q\). The equality of extension
classes gives
\((1+\varepsilon t)(b+\varepsilon\beta)=b(1+\varepsilon\bar r)\), or
\(\beta=b\bar r-tb\). Conversely such a pair gives an actual
isomorphism of extensions and a trivial module deformation.
Every matrix over \(A\) lifts to \(R\). Therefore

\[
 \ker L=\operatorname{im}\psi,\qquad
 e_1:=\operatorname{ext}^1_R(M,M)
       \ge q^2(2a-2\ell)-(q^2\ell+c-a_L)
       =q^2(2a-3\ell)-c+a_L.
 \tag{NI.14}
\]

No assertion that this fixed-ideal family exhausts deformations
of \(M\) is needed. It supplies the stated lower bound.

\subsection{4. A general finite-hull-cone Euler identity}\label{SF-E849:s4}

We prove the cone identity in greater generality to display
exactly where equal lengths enter. Let \(N\) have a canonical
locally free hull \(R^h\), finite kernel \(T_N\), and finite quotient \(Q_N\).
Put

\[
 \begin{split}
 \tau&=\operatorname{length}T_N,&
 \zeta&=\operatorname{length}Q_N,&
 C&=\operatorname{cone}(N\longrightarrow R^h),\\
 k_N&=\dim\operatorname{Hom}_R(N,T_N),&
 d_N&=\dim\operatorname{Hom}_R(Q_N,T_N),\\
 \mathcal B_N&=\operatorname{im}(\operatorname{End}_R(N)
                         \longrightarrow\operatorname{Mat}_h(R)),&
 c_N&=\operatorname{length}(\operatorname{Mat}_h(R)/\mathcal B_N).
 \end{split}
 \tag{NI.15}
\]

Let \(J\) be the hull image. Matrices whose columns lie in \(J\) belong
to \(\mathcal B_N\): lift each column through \(N\to J\) and compose the resulting
map \(R^h\to N\) with the hull map. This column module has quotient
of length \(h\zeta\) in \(\operatorname{Mat}_h(R)\), so \(c_N\) is finite. The kernel of
\(\operatorname{End}(N)\to \mathcal B_N\) is \(\operatorname{Hom}(N,T_N)\).

The grade-four finite-defect vanishings give
\(\operatorname{Hom}_R(N,R^h)=\operatorname{Mat}_h(R)\) and
\(\operatorname{Ext}^1_R(N,R^h)=\operatorname{Ext}^2_R(N,R^h)=0\). Write
\(e_i=\operatorname{ext}^i_R(N,N)\), and \(h_j=\operatorname{ext}^j_R(N,C)\).
The hull triangle immediately gives
\(h_{-1}=k_N\), \(h_0=e_1+c_N\), and \(h_1=e_2\).

All cohomology of \(F=\operatorname{RHom}_R(N,C)\) has finite length. Supported
duality and perfect adjunction give
\(\operatorname{Ext}^j(N,C)^\vee=\operatorname{Ext}^{4-j}(C,N)\). Applying \(\operatorname{Hom}_R(-,N)\) to the
hull triangle gives \(\operatorname{Ext}^2(C,N)=\operatorname{Ext}^1(N,N)\), so \(h_2=e_1\).
It also gives \(\operatorname{Hom}(C,N)=\operatorname{Hom}(Q_N,N)=\operatorname{Hom}(Q_N,T_N)\), of dimension
\(d_N\); the finite image of \(Q_N\) must lie in the canonical torsion.
Thus \(h_4=d_N\).

To compute \(h_3\), consider the image of \(\operatorname{Hom}(R^h,N)\) in \(\operatorname{End}(N)\).
Its hull image is exactly the column module \(J^h\). The intersection
of that image with \(\operatorname{Hom}(N,T_N)\) is the image of \(\operatorname{Hom}(R^h,T_N)\);
the latter map has kernel \(\operatorname{Hom}(Q_N,T_N)\). Its dimension is
\(h\tau-d_N\). The quotient \(\mathcal B_N/J^h\) has dimension \(h\zeta-c_N\).
The exact cokernel description of \(\operatorname{Ext}^1(C,N)\) therefore gives

\[
 \begin{split}
 (h_{-1},h_0,h_1,h_2,h_3,h_4)
  =\bigl(k_N,\ e_1+c_N,\ e_2,\ e_1,\
        k_N+d_N+h(\zeta-\tau)-c_N,\ d_N\bigr).
 \end{split}
 \tag{NI.16}
\]

For the asserted intersection, an image endomorphism has zero
hull matrix exactly when each lifted column is in \(T_N\).
Conversely every such column map gives an endomorphism in that
intersection. Its kernel consists of maps factoring through the
quotient \(R^h/J=Q_N\). This proves the dimension formula without
assuming that \(T_N\) is annihilated by \(J\).

A finite free resolution of \(N\) has alternating ranks summing
to \(h\). Since C has finite Euler length \(\zeta-\tau\), it follows that

\[
 \chi_{\rm length}R\operatorname{Hom}_R(N,C)
       =h(\zeta-\tau),\qquad
 e_2=2e_1-2k_N+2c_N+2h(\tau-\zeta).
 \tag{NI.17}
\]

The second identity is obtained by substituting all six finite
dimensions (NI.16) with their signs into the first. There is no
subtraction of infinite-dimensional End groups and no assumed
ordinary self-duality of \(\operatorname{RHom}(N,N)\).

For \(M_b\) in (NI.2), the two defect lengths are equal, \(h=q\), and
\(k_N=q^2a+a_L\). Combining (NI.14) and (NI.17) yields

\[
 \begin{split}
 \operatorname{ext}^2_R(M_b,M_b)
  &\ge2\bigl(q^2(2a-3\ell)-c+a_L\bigr)
           -2(q^2a+a_L)+2c\\
  &=2q^2(a-3\ell)=2q^2\delta.
 \end{split}
 \tag{NI.18}
\]

This proves the local theorem. The hull-image correction and
the left annihilator have cancelled without any CI assumption.

\subsection{5. Actual finite stabilizers and the rectangular-to-square step}\label{SF-E849:s5}

Suppose (NI.4)--(NI.5) and the evaluation hypothesis of E849b.
The stabilizer of \(p\) under the translation/tensor action has
identity translation: a translation fixing a point is identity.
Torsion Picard factors admit flat local frames, with one constant
projective scalar cocycle \(\alpha\). Thus the actual completed module
has a decomposition over the semisimple twisted finite group
algebra

\[
 \widehat E_p=\bigoplus_\lambda S_\lambda\otimes N_\lambda.
 \tag{NI.19}
\]

The group acts on the simple representation factor \(S_{\lambda}\);
it acts \(R\)-linearly and trivially on the base coordinates.
This is the exact point-stabilizer interface of E846--E848 and
supplied S21. It is not a claim for arbitrary finite formal
automorphisms. Maschke averaging is valid for the projective
cocycle because the scalar factors cancel in conjugation.

Apply the exact multiplicity functor to the canonical torsion,
hull, image, and quotient. The torsion and quotient multiplicities
are direct summands of finite direct sums of \(A^m\), since the
averaging idempotents are A-linear. They are free over the local
algebra \(A\); write their ranks as \(t_{\lambda}\) and \(q_{\lambda}\). The hull
multiplicity is \(R\)-free, of rank \(h_{\lambda}\). A surjective matrix
\(R^{h_{\lambda}}\to A^{q_{\lambda}}\) has a unit minor modulo the maximal
ideal. Lifting elementary operations puts it into the ordinary
residue map on \(q_{\lambda}\) coordinates and zero on the others.
Its kernel is \(I^{q_{\lambda}}\) direct-sum \(R^{h_{\lambda}-q_{\lambda}}\).
The extension over the free summand splits, giving

\[
 N_\lambda\simeq M_\lambda\oplus R^{h_\lambda-q_\lambda},
 \qquad
 0\longrightarrow A^{t_\lambda}\longrightarrow M_\lambda
       \longrightarrow I^{q_\lambda}\longrightarrow0.
 \tag{NI.20}
\]

This normal form requires only that \(A=R/I\) be local, not that \(I\)
be generated by a regular sequence.

Ordinary evaluation is natural for restriction, completion, and
the actual direct sum (NI.19). The \(H^1(T_X)\) and \(H^2(\mathcal O_X)\) HKR
components vanish on the formal affine neighborhood. In abelian
logarithm coordinates the remaining bivector \(\pi\) is a nonzero
constant tensor. Hence \(\rho_{\pi}(N_{\lambda})=0\) for each \(\lambda\).

The exact E819 local implication applies to each of these
ordinary multiplicity modules. Vanishing of evaluation permits
a continuous first-order free-resolution lift over the algebra
\(a*b=ab+\varepsilon\pi(da,db)/2\). Grade-four vanishing
\(\operatorname{Ext}^1_R(N_{\lambda},R^{h_{\lambda}})=0\) lifts its hull map. The cone
has finite cohomology and a bounded finite free representative
over the dual numbers, by the finite Tor amplitude criterion
in Stacks Lemma 15.76.2, tag 0658.\cite{SF-StacksPerfect} The coordinate commutator on
this perfect cone has zero supertrace; for any nonzero \(\pi_{ij}\)
this says

\[
 0=\varepsilon\pi_{ij}
       \bigl(\operatorname{length}Q_\lambda
                  -\operatorname{length}T_\lambda\bigr)
   =\varepsilon\pi_{ij}(q_\lambda-t_\lambda)\ell.
 \tag{NI.21}
\]

Therefore \(t_{\lambda}=q_{\lambda}\). This argument did not use the CI
property; it is the finite-cone trace proof of E819, retained
with precisely its first-order scope. If \(q_{\lambda}=0\) it also
forces \(t_{\lambda}=0\). Some \(q_{\lambda}\) is positive because \(A^m\) is
nonzero. For that component the square theorem (NI.3) applies.
No equivariant noncommutative deformation and no chosen-basis
conjugacy of the whole action has been assumed.

\subsection{6. Global support, trace, and the invariant kernel}\label{SF-E849:s6}

Let \(C_X=\operatorname{cone}(E\to H)\). Its cohomology has finite support \(Z\). The
cohomology sheaves of \(\operatorname{RHom}(E,C_X)\) are finite-supported, so their
higher cohomology vanishes. Restriction and completion give

\[
 \operatorname{Ext}^j_X(E,C_X)
       \simeq\bigoplus_{p'\in Z}
          \operatorname{Ext}^j_{\widehat{\mathcal O}_{X,p'}}
                             (\widehat E_{p'},C_{p'}).
 \tag{NI.22}
\]

The local grade-four vanishings identify the \(j=1\) group on the
right with local \(\operatorname{Ext}^2(E,E)\). The preceding global arrow
\(\operatorname{Ext}^1_X(E,H)\to \operatorname{Ext}^1_X(E,C_X)\) is zero: in its natural restriction
square the bottom source is zero and the right vertical arrow
is the isomorphism (NI.22). The global hull triangle therefore
gives the actual canonical injection

\[
 j:\bigoplus_{p'\in Z}
        \operatorname{Ext}^2_{\widehat{\mathcal O}_{X,p'}}
                        (\widehat E_{p'},\widehat E_{p'})
                 \hookrightarrow\operatorname{Ext}^2_X(E,E).
 \tag{NI.23}
\]

This argument uses the finite-support isomorphism to prove the
global arrow vanishes. It does not infer vanishing of a global
connecting map solely from a vanishing localized source.

The Serre pairing of \(j(u)\) with \(\rho_E(\gamma)\), for \(\gamma\) in \(HH^2\),
is the trace of the connecting image of a \(\rho_E(\gamma)\), where
a in \(\operatorname{Ext}^1(E,C_X)\) corresponds to \(u\). That product lies in
\(\operatorname{Ext}^3(E,C_X)\), so by (NI.22) it only depends on the completed
restriction of the evaluation. The higher HKR components
vanish there; only the six-dimensional space
\(H^0(X,\bigwedge^2T_X)\) can contribute. The full CY4 trace/evaluation
adjunction therefore bounds the rank of semiregularity on
every subspace of (NI.23) by six.

The exact adjunction is the supplied DEG\_L S20 first proof
paragraph and S21.6 through the adjunction preceding S21.8.
The ordinary global Serre pairing is Stacks 48.27,
Lemma 48.27.1(5),(7), Remarks 48.27.2--3, and Lemma 48.27.4.\cite{SFStacksE846Duality}
This uses the full adjunction, not just the composition
formula \(\sigma\rho\) equals contraction by the character.

Choose \(q_{\lambda}>0\) as in Section 5. The classes \(\operatorname{Id}_{S_\lambda}\otimes u\), for \(u\) in \(\operatorname{Ext}^2_R(M_{\lambda},M_{\lambda})\), are stabilizer
invariant and injective as a vector space. Place them in their
actual local block and zero in the other blocks; transport and
sum over the orbit of p. Stabilizer invariance makes the sum
well-defined, and its component at \(p\) proves injectivity. The
canonical \(j\) takes this space into invariant global Ext. The
six-dimensional trace bound persists over the orbit. Supplied
S21 identifies this ordinary invariant map with the relevant
identity-sector semiregularity map for the effective action.
Consequently

\[
 \dim\ker\sigma^2_{E,G}
       \ge2q_\lambda^2\delta-6
       \ge2\delta-6\ge2.
 \tag{NI.24}
\]

This proves E849b. For the identity action the given square
block itself supplies this space, so the evaluation hypothesis
and Section 5 are not needed.

For (NI.6), the two balanced-plane cap-kernel classes have
nonzero bivectors: \(\pi_++\nu\eta_+\) and \(\pi_-+\nu\eta_-\), where
\(\eta_{\pm}=\pi_{\pm}\) contracted with \(\theta_{\pm}^2/2\). The two terms cancel
on the entire corresponding balanced factor. Conjugation by
the algebraic B-field preserves the bivector part. Published
Buchweitz--Flenner Theorem 6.4.2 and Proposition 6.2.1\cite{SF-BF2008} give
\(\sigma\rho\) equals contraction by the entire character, and supplied S21
gives its equivariant compatibility. Hypothetical injectivity
of \(\sigma_G\) would force the ordinary evaluation of one of these
classes to vanish. E849b then gives the contradiction (NI.24).
This proves E849c, with the full B-field and coherent action.

\subsection{7. The square-zero non-CI example and strict scope}\label{SF-E849:s7}

Take \(I=\mathfrak m^2\). The algebra \(A\) has basis \(1,x_1,x_2,x_3,x_4\), with
all quadratic products zero. Its socle is the span of the four
variables, so its dimension is four. An Artinian complete
intersection has one-dimensional socle, so this algebra is not
one. This also follows from its ten independent minimal
quadratic generators in a regular local ring of dimension four.

We compute \(\operatorname{Hom}_R(I,A)\) without a Hilbert-scheme dimension claim.
The ten degree-two monomials give a basis of \(I/\mathfrak m I\). Every
homomorphism \(I\to A\) has zero constant term on each generator.
Indeed for a generator \(x_ix_j\) choose \(k\) different from \(i\).
The relation
\(x_k(x_ix_j)-x_i(x_kx_j)=0\) implies
\(c_{ij}x_k-c_{kj}x_i=0\) in \(A\), where c denotes the constant
coefficient of an image. The two linear basis vectors are
independent, so \(c_{ij}=0\).

Conversely every complex-linear map \(I/\mathfrak m I\to\mathfrak m_A\) gives an
\(R\)-linear map \(I\to A\): all positive-degree elements act as zero
on both \(I/\mathfrak m I\) and \(\mathfrak m_A\). Thus

\[
 \operatorname{Hom}_R(\mathfrak m^2,R/\mathfrak m^2)
     \simeq\operatorname{Hom}_{\mathbf C}
           (\operatorname{Sym}^2\mathbf C^4,\mathbf C^4),
 \qquad a=10\cdot4=40.
 \tag{NI.25}
\]

Equations (NI.1), (NI.3), and (NI.10) now give \(\ell=5\),
\(\delta=25\), \(\dim V=70\), and (NI.7). In the global balanced
application the invariant kernel has dimension at least
\(50q_{\lambda}^2-6\), in particular at least forty-four.

There are actual coherent examples with these local modules.
At a point of the fixed CM RM fourfold, take the algebraic
subscheme \(Z\) defined by \(\mathfrak m_p^2\). Let \(H_0\) be the actual rank-32
balanced vector bundle of E815/E820, and take \(H_0^{\oplus q}\).
A local frame gives a surjection of this hull to \(\mathcal O_Z^q\).
Its kernel has completed type \(I^q\) direct-sum \(R^{31q}\).
Every class in \(\operatorname{Mat}_q(V)\) is realized by an actual extension
of that kernel by \(\mathcal O_Z^q\): the target has finite support, so the
Ext spectral sequence identifies global Ext-one with the
completed local group. The resulting sheaf has locally free
hull \(H_0^{\oplus q}\), equal torsion and hull quotient \(\mathcal O_Z^q\),
and full character equal to that of the hull. It is ordinarily
gluable because it is coherent. The identity group is an
actual action, and the theorem addresses every additional
specified effective coherent action if one exists.

This example is outside E836--E848's common complete-intersection
scope, outside the reduced and curvilinear cases, and has equal
signed point-defect length, so E819's length equality alone
does not rule it out. It supplies a strict new local-algebra
case for the auxiliary fourfold programme. No successful
semiregular object is produced.

\subsection{8. Dependencies and remaining boundaries}\label{SF-E849:s8}

The local fixed-ideal tangent family, its exact orbit kernel,
the general cone identity with unequal lengths, and the \(\mathfrak m^2\)
calculation are proved in full here. The square CI precursor
E848 had joint root/sixfolds provenance; no independent origin
for that precursor is claimed. The new proof does not alter
its frozen text.

The published local/global duality and perfectness inputs are
the precisely located inputs already checked in E847--E848.
The ordinary first-order module/hull/cone argument is accepted
E819, with no dependence on a CI ideal; its exact hypothesis
is repeated in Section 5. The BF input and supplied S20/S21
keep their published and supplied statuses respectively.
No Markman or unrefereed deformation theorem is invoked in
the new local bound.

The result still requires a common cyclic Artin algebra with
torsion and quotient free over that algebra, and the numerical
condition \(\delta\ge4\) for its uniform finite-action conclusion.
Different Artin algebras, arbitrary nonfree finite modules,
nonfree hulls, positive-dimensional defects, nontrivial formal
base automorphisms, and general derived objects are not
included. The fixed-ideal family does not establish a tangent
bound for arbitrary varying quotient ideals.
